\documentclass[10pt,a4paper]{amsart}
\usepackage[margin=19mm]{geometry}
\usepackage{amsmath,amssymb,mathtools}
\usepackage[scr=rsfs,cal=euler]{mathalfa}
\usepackage{xfrac}
\usepackage[unicode,hidelinks,bookmarks=false]{hyperref}
\newcommand{\ie}{i.e.\ }
\newcommand{\cf}{cf.\ }
\newcommand{\resp}{respectively\ }
\newcommand{\Subsection}{\subsection}
\providecommand{\nobreakdash}{\nobreak\hbox{-}}
\setlength{\emergencystretch}{2em}
\multlinegap0pt
\usepackage{mathtools}
\usepackage{amssymb}
\usepackage{float}
\usepackage{enumerate}
\usepackage{tikz}
\usepackage[shortlabels]{enumitem}
\newtheorem{lem}{Lemma}
\def\dd{{\rm d}}
\title{The relativistic Euler-Poisson equation as a mean-field limit}
\author{Immanuel Ben-Porat}
\date{Source: arXiv:2609.14880v1 (CC BY 4.0)}
\begin{document}
\maketitle

\noindent\emph{Source and licence.} The relativistic Euler-Poisson equation as a mean-field limit. Version arXiv:2609.14880v1 (CC BY 4.0), distributed by arXiv under the Creative Commons Attribution 4.0 International licence, http://creativecommons.org/licenses/by/4.0/, as recorded in the arXiv licence metadata for this exact version and shown on the abstract page https://arxiv.org/abs/2609.14880v1. Full licence notice: \texttt{licenses/arxiv-2609.14880-CC-BY-4.0.txt}.\par\smallskip

\noindent\textit{Editorial scope.} Counted unit: the article's lem eta (source0.tex lines 429--435). The article's complete original proof of it is delivered with it (source0.tex lines 436--536, one \texttt{proof} environment of 101 lines). No mathematics is written, repaired or re-derived: the statement and the proof are the article's own lines, in the article's own notation, and no retained line is substituted or re-flowed.  Retained uncounted as the source-local context the proof needs: lines 404--409.  Nothing was omitted from any retained span, so the package carries no omission record.  No retained line contains a citation, so the wrapper carries no bibliography and no bibliography entry.  No retained line contains a figure environment, an image-inclusion command or drawing code, so no image is delivered and the package declares no asset dependency.  Editorial formatting only, in addition: an AMS \texttt{amsart} preamble that restates the article's own declarations and macros used by the retained lines, namely the environments \texttt{lem} on separate counters, together with the standard packages those lines need.  The retained environments print in this document as Lemma 1 in order of appearance, whereas the article prints the same items with the numbering of its own full document.  The complete native source of the article as distributed in the arXiv e-print for this exact version is bundled unchanged as \texttt{sources/210368/source0.tex}.
\begin{lem}\label{elementary ine about min}
For all $r\geq 0$ it holds that 
\begin{align*}
\frac{r^{2}}{1+r}\leq \min\{r,r^{2}\}\leq \frac{2r^{2}}{1+r}.     
\end{align*}
\end{lem}

\begin{lem}\label{coercivity positivity of eta}
There is a constant $C_{R}>0$ such that for any $(p_{1},p_{2})\in \mathbb{R}^{d}\times B_{R}(0)$ it holds that 
 \begin{align*}
 \eta(p_{1}\vert p_{2})\geq \frac{C_{R}\left\vert p_{1}-p_{2}\right\vert^{2}}{1+\left\vert p_{1}-p_{2}\right\vert}.     
 \end{align*}
In particular, $\eta(p_{1}\vert p_{2})\geq 0$ for all $p_{1},p_{2}\in \mathbb{R}^{d}$. 
\end{lem}

\begin{proof}
In view of Lemma \ref{elementary ine about min} it suffices to prove that there is some $C_{R}>0$ such that for all $(p_{1},p_{2})\in \mathbb{R}^{d}\times B_{R}(0)$ it holds that 
\begin{align}
\eta(p_{1}\vert p_{2})\geq C_{R}\min\{\left\vert p_{1}-p_{2}\right\vert ,\left\vert p_{1}-p_{2}\right\vert^{2} \}. \label{eta requested inequality}   \end{align}
To achieve this inequality we proceed via the following steps. \\
\textbf{Calculation of the eigenvalues of $\nabla^{2}\mathscr{K}(z)$.} We claim that for any $z\neq 0$ the Hessian matrix $\nabla^{2}\mathscr{K}(z)$ has exactly $2$ positive eigenvalues (up to multiplicity) given by 
\begin{align}
\lambda_{\max}(z)=\frac{1}{\sqrt{1+\left\vert z\right\vert^{2}}} \ \mbox{and}\ \lambda_{\min}(z)=\frac{1}{(1+\left\vert z\right\vert^{2})^{\frac{3}{2}}}. \label{formula for eigenvalues}    
\end{align}
First, by direct calculation we find that the Hessian of $\mathscr{K}$ is given by 
\begin{align*}
\nabla^{2}\mathscr{K}(z)=\frac{1}{\sqrt{1+\left\vert z\right\vert^{2} }}I-\frac{z\otimes z}{(1+\left\vert z\right\vert^{2})^{\frac{3}{2}}}.     
\end{align*}
If $v\perp z$ then we have 
\begin{align*}
(z\otimes z)v=(z\cdot v)z=0.   
\end{align*}
Therefore it follows that 
\begin{align*}
\nabla^{2}\mathscr{K}(z)v=\frac{1}{\sqrt{1+\left\vert z\right\vert^{2}}}Iv=\frac{1}{\sqrt{1+\left\vert z\right\vert^{2}}}v.
\end{align*}
It follows that $\frac{1}{\sqrt{1+\left\vert z\right\vert^{2}}}$ is an eigenvalue with eigenspace $z^{\perp}$. Therefore, since $\nabla^{2}\mathscr{K}(z)$ is symmetric the multiplicity of $\frac{1}{\sqrt{1+\left\vert z\right\vert^{2}}}$ is $\dim(z^{\perp})=d-1$. On the other hand, using the identity $$(z\otimes z)z=\left\vert z\right\vert^{2}z$$ we get  
\begin{align*}
\nabla^{2}\mathscr{K}(z)z=\frac{1}{\sqrt{1+\left\vert z\right\vert^{2} }}z-\frac{\left\vert z\right\vert^{2} }{(1+\left\vert z\right\vert^{2} )^{\frac{3}{2}}}z=\left(\frac{1}{\sqrt{1+\left\vert z\right\vert^{2}}}-\frac{\left\vert z\right\vert^{2}}{(1+\left\vert z\right\vert^{2} )^{\frac{3}{2}}}\right)z=\frac{1}{(1+\left\vert z\right\vert^{2})^{\frac{3}{2}}}z.    
\end{align*}
It follows that $\frac{1}{(1+\left\vert z\right\vert^{2})^{\frac{3}{2}}}$ is an eigenvalue with eigenspace $\mathrm{span}\{z\}$. Since $z\neq 0$ both of the eigenvalues we found are distinct. Moreover $\dim(z^{\perp})=d-1$ while $\dim(\mathrm{span}\{z\}{})=1$ and therefore these are the only eigenvalues. 
To conclude, for any $z\neq 0$ the Hessian  $\nabla^{2}\mathscr{K}(z)$ has exactly two distinct eigenvalues given by \ref{formula for eigenvalues}.
We will now separate between two regions: Fix $M=M_{R}$ to be chosen later.  We will first consider the case where $\left\vert p_{1}-p_{2}\right\vert\leq M $ and then the case where $\left\vert p_{1}-p
_{2}\right\vert\geq M$. We may assume with no loss of generality that $p_{1}\neq p_{2}$. \\ 
\textbf{The case where $\left\vert p_{1}-p_{2}\right\vert\leq M$}. By the theorem of Taylor-Hadamard we have 
\begin{align*}
\eta(p_{1}\vert p_{2})&=\mathscr{K}(p_{1})-\mathscr{K}(p_{2})-\nabla \mathscr{K}(p_{2})\cdot (p_{1}-p_{2})\\
&=\int_{0}^{1}(\nabla \mathscr{K}(p_{2}+s(p_{1}-p_{2}))-\nabla \mathscr{K}(p_{2}))\cdot (p_{1}-p_{2})\ \dd s\\
&= \int_{0}^{1} s(p_{1}-p_{2})^{T}\nabla^{2}\mathscr{K}(\xi_{s})(p_{1}-p_{2})\ \dd s
\end{align*}
where in the last equation we applied the intermediate value theorem with an intermediate point $\xi_{s}=(1-\theta )p_{2}+\theta(p_{2}+s(p_{1}-p_{2}))=(1-\theta s)p_{2}+\theta sp_{1}$ for some $\theta \in [0,1]$. Consequently, we deduce that \footnote{Recall that if $A$ is a matrix with minimal eigenvalue $\lambda_{\min}$ then $\lambda_{\min}=\underset{\left\vert v\right\vert=1 }{\min}v^{T}Av$. }  
\begin{align}
\eta(p_{1}\vert p_{2})\geq \left\vert p_{1}-p_{2}\right\vert^{2} \int_{0}^{1}s\min_{\left\vert v\right\vert=1 }v^{T}\nabla^{2}\mathscr{K}(\xi_{s})v\ \dd s
=\left\vert p_{1}-p_{2}\right\vert^{2} \int_{0}^{1}s\lambda_{\min}(\xi_{s})\ \dd s. \label{eta(p1p2) est in term integral}   
\end{align}
In view of \eqref{formula for eigenvalues} we have that 
\begin{align*}
\int_{0}^{1}s\lambda_{\min}(\xi_{s})\ \dd s&=\int_{0}^{1}\frac{s}{(1+\left\vert \xi_{s}\right\vert^{2})^{\frac{3}{2}}}\ \dd s\\
&=\int_{0}^{1}\frac{s}{(1+\left\vert (1-\theta s)p_{2}+\theta s p_{1}\right\vert^{2})^{\frac{3}{2}}}\ \dd s.    \end{align*}
In addition, note that  
\begin{align*}
\left\vert (1-\theta s)p_{2}+\theta sp_{1}\right\vert^{2}&\leq 2(1-\theta s)^{2}\left\vert p_{2}\right\vert^{2}+2\theta^{2}s^{2} \left\vert p_{1}\right\vert^{2}\\
&\leq 2\left\vert p_{2}\right\vert^{2}+4(\left\vert p_{1}-p_{2}\right\vert^{2}+\left\vert p_{2}\right\vert^{2} 
)\leq 6R^{2}+4M^{2}.     \end{align*}
Hence, we infer  
\begin{align}
\int_{0}^{1} s\lambda_{\min}(\xi_{s})\ \dd s\geq \frac{1}{(1+6R^{2}+4M^{2})^{\frac{3}{2}}}. \label{lower bound on integral}    
\end{align}
Substituting \eqref{lower bound on integral} inside \eqref{eta(p1p2) est in term integral} we conclude that 
\begin{align}
\eta(p_{1}\vert p_{2})\geq \frac{1}{(1+6R^{2}+4M^{2})^{\frac{3}{2}}}\left\vert p_{1}-p_{2}\right\vert^{2}. \label{quadratic est on eta}   
\end{align}
\textbf{The case where $\left\vert p_{1}-p_{2}\right\vert\geq M $.} First, we claim to have the inequality 
\begin{align}
\mathscr{K}(p_{1})-\nabla \mathscr{K}(p_{2})\cdot p_{1}\geq \left\vert p_{1}\right\vert\left(1-\frac{R}{\sqrt{1+R^{2}}}\right).\label{inequality for K}     
\end{align}
Indeed 
\begin{align*}
\mathscr{K}(p_{1})-\nabla \mathscr{K}(p_{2})\cdot p_{1}&=\sqrt{1+\left\vert p_{1}\right\vert^{2}}-\frac{p_{2}\cdot p_{1}}{\sqrt{1+\left\vert p_{2}\right\vert^{2}  }}\\
&\geq \left\vert p_{1}\right\vert-\frac{\left\vert p_{1}\right\vert\left\vert p_{2}\right\vert  }{\sqrt{1+\left\vert p_{2}\right\vert^{2}}}=\left\vert p_{1}\right\vert\left(1-\frac{\left\vert p_{2}\right\vert}{\sqrt{1+\left\vert p_{2}\right\vert^{2} }}\right).       
\end{align*}
Since $r\mapsto \frac{r}{\sqrt{1+r^{2}}}$ is increasing and $\left\vert p_{2}\right\vert\leq R$ it follows that 
\begin{align*}
1-\frac{\left\vert p_{2}\right\vert}{\sqrt{1+\left\vert p_{2}\right\vert^{2}}}\geq 1-\frac{R}{\sqrt{1+R^{2}}}    
\end{align*}
which establishes \eqref{inequality for K}. 
Owing to \eqref{inequality for K} and noticing the identity 
$$\nabla \mathscr{K}(p_{2})\cdot p_{2}-\mathscr{K}(p_{2})=\frac{\left\vert p_{2}\right\vert^{2} }{\sqrt{1+\left\vert p_{2}\right\vert^{2} }}-\sqrt{1+\left\vert p_{2} \right\vert^{2} }=-\frac{1}{\sqrt{1+\left\vert p_{2}\right\vert^{2}}}$$we infer that 
\begin{align*}
\eta(p_{1}\vert p_{2})&=
\mathscr{K}(p_{1})-\nabla \mathscr{K}(p_{2})\cdot p_{1}+\nabla \mathscr{K}(p_{2})\cdot p_{2}-\mathscr{K}(p_{2})\\
&\geq \left\vert p_{1}\right\vert\left(1-\frac{R}{\sqrt{1+R^{2}}}\right)-\frac{1}{\sqrt{1+\left\vert p_{2}\right\vert^{2}}}\\
&\geq \left\vert \left\vert p_{2}-p_{1}\right\vert-\left\vert p_{2}\right\vert\right\vert\left(1-\frac{R}{\sqrt{1+R^{2}}}\right)-\frac{1}{\sqrt{1+\left\vert p_{2}\right\vert^{2}}}. \end{align*}
Since $\left\vert p_{2}\right\vert\leq R $ and $\left\vert p_{1}-p_{2}\right\vert\geq M$, we have $\left\vert \left\vert p_{1}-p_{2}\right\vert -\left\vert p_{2}\right\vert \right\vert=\left\vert p_{1}-p_{2}\right\vert  -\left\vert p_{2}\right\vert $ provided $M\geq R$.  So, as long as $M\geq R$ we get 
\begin{align*}
\eta(p_{1}\vert p_{2})&\geq \left\vert p_{1}-p_{2}\right\vert\left(1-\frac{R}{\sqrt{1+R^{2}}}\right)-\left\vert p_{2}\right\vert\left(1-\frac{R}{\sqrt{1+R^{2}}}\right)-\frac{1}{\sqrt{1+\left\vert p_{2}\right\vert^{2}}}\\
&\geq \left\vert p_{1}-p_{2}\right\vert\left(1-\frac{R}{\sqrt{1+R^{2}}}\right)-R\left(1-\frac{R}{\sqrt{1+R^{2}}}\right)-1. 
\end{align*}
Setting 
\begin{align*}
\alpha_{R}\coloneqq 1-\frac{R}{\sqrt{1+R^{2}}}\ \mbox{and}\ \beta_{R}\coloneqq R\left(1-\frac{R}{\sqrt{1+R^{2}}}\right)+1    
\end{align*}
the latter inequality writes 
\begin{align*}
\eta(p_{1}\vert p_{2})\geq \alpha_{R}\left\vert p_{1}-p_{2}\right\vert-\beta_{R}.      
\end{align*}
Now set $M\coloneqq \max\{R,\frac{2\beta_{R}}{\alpha_{R}}\}$ and note that 
\begin{align*}
\alpha_{R}\left\vert p_{1}-p_{2}\right\vert-\beta_{R}\geq \frac{\alpha_{R}}{2}\left\vert p_{1}-p_{2}\right\vert\iff \left\vert p_{1}-p_{2}\right\vert\geq \frac{2\beta_{R}}{\alpha_{R}}.        
\end{align*}
So choosing  $M=\max\{R,\frac{2\beta_{R}}{\alpha_{R}}\}$ we conclude that 
\begin{align}
\eta(p_{1}\vert p_{2})\geq \frac{\alpha_{R}}{2}\left\vert p_{1}-p_{2}\right\vert. \label{linear lower bound}     
\end{align}
Combining \eqref{quadratic est on eta} with \eqref{linear lower bound} yields \eqref{eta requested inequality}. 
\end{proof}

\end{document}
